% ==STYLE-BEGIN==
\documentclass[UTF8,11pt]{ctexart}
\usepackage{amsmath,amssymb,amsthm}
\usepackage[a4paper,margin=2.4cm,top=2.3cm,bottom=2.3cm]{geometry}
\usepackage{fancyhdr}
\usepackage{enumitem}
\usepackage{hyperref}
\hypersetup{hidelinks}

% ---- 版式 ----
\linespread{1.15}
\setlength{\parskip}{0.35em}
\setlength{\parindent}{1.6em}
\setlist{nosep,leftmargin=2em}

% ---- 页眉 ----
\pagestyle{fancy}
\fancyhf{}
\fancyhead[L]{\small\coursename}
\fancyhead[R]{\small\thepage}
\renewcommand{\headrulewidth}{0.4pt}

% ---- 定理环境：中文用正体，标题加粗，按小节共用一条流水号 ----
\newtheoremstyle{cnthm}{0.7em}{0.7em}{\normalfont}{}{\bfseries}{}{0.5em}{}
\theoremstyle{cnthm}
\newtheorem{thm}{定理}[section]
\newtheorem{lem}[thm]{引理}
\newtheorem{prop}[thm]{命题}
\newtheorem{cor}[thm]{推论}
\newtheorem{defn}[thm]{定义}
\newtheorem{rem}[thm]{注}
\newtheorem{exa}[thm]{例}
\newtheorem{conj}[thm]{猜想}
\renewcommand{\proofname}{\bfseries 证明}
\providecommand{\coursename}{}
% ==STYLE-END==
\renewcommand{\coursename}{高等代数\,·\,1026}
\title{高等代数 1026}
\date{}
\begin{document}
\maketitle

\section{一元 $n$ 次方程与多项式}

\begin{defn}[一元 $n$ 次方程]
设 $K$ 是数域，$a_0,\cdots,a_n\in K$，且 $a_0\neq0$，$n\geq1$，称
\[
a_0x^n+a_1x^{n-1}+\cdots+a_n=0
\]
为 $K$ 上的一元 $n$ 次方程。若 $x=\alpha$ 满足方程，则称 $\alpha$ 是方程的解或根。
\end{defn}

\begin{thm}[代数基本定理]
上述方程在 $\mathbb{C}$ 上一定有解；等价地，$\forall\deg f(x)\geq1$ 的 $f(x)$ 在 $\mathbb{C}$ 上必有零点。
\end{thm}

\begin{defn}[一元 $n$ 次多项式]
称 $a_0x^n+\cdots+a_n$（$a_0\neq0$）是 $K$ 上的一元 $n$ 次多项式，用 $f(x),g(x)$ 等表示。$a_0$ 称为首项系数，$n$ 称为次数，用 $\deg f(x)$ 表示 $f(x)$ 的次数，例如
\[
\deg(c)=0,\qquad \deg(0)=-\infty.
\]
若 $\alpha\in\mathbb{C}$ 满足 $f(\alpha)=0$，则称 $\alpha$ 是 $f(x)$ 的零点。
\end{defn}

\section{因式分解与零点}

\begin{lem}[除以 $x-\alpha$ 的带余除法]
若 $f(x)$ 是 $K$ 上的 $n$ 次多项式，$\forall\alpha\in\mathbb{C}$，则存在首项系数为 $a_0$ 的 $(n-1)$ 次多项式 $g(x)$，使得
\[
f(x)=(x-\alpha)g(x)+f(\alpha).
\]
\end{lem}

\begin{proof}
\[
\begin{aligned}
f(x)-f(\alpha)
&=a_0(x^{n}-\alpha^{n})+a_1(x^{n-1}-\alpha^{n-1})+\cdots+a_{n-1}(x-\alpha)\\
&=a_0(x-\alpha)(x^{n-1}+x^{n-2}\alpha+\cdots+\alpha^{n-1})
  +a_1(x-\alpha)(\cdots)+\cdots\\
&=(x-\alpha)\bigl[\,a_0(x^{n-1}+x^{n-2}\alpha+\cdots+\alpha^{n-1})+a_1(\cdots)+\cdots\bigr].
\end{aligned}
\]
这里用了 $x^{k}-\alpha^{k}=(x-\alpha)(x^{k-1}+x^{k-2}\alpha+\cdots+\alpha^{k-1})$（$1\leq k\leq n$）。

令 $g(x)$ 为上式方括号内的多项式。由于 $x^{n-1}$ 这一项只能由第一项
$a_0(x^{n-1}+x^{n-2}\alpha+\cdots+\alpha^{n-1})$ 产生，后面每个 $a_i$ 对应的括号最高次都低于 $n-1$，故 $g(x)$ 的最高次项为 $a_0x^{n-1}$，即 $g(x)$ 是首项系数为 $a_0$ 的 $(n-1)$ 次多项式，且 $f(x)=(x-\alpha)g(x)+f(\alpha)$。
\end{proof}

特别地，若 $\alpha$ 是 $f(x)$ 的零点，则 $f(x)=(x-\alpha)g(x)$。

\begin{prop}
若 $f(x)$ 是首项系数为 $a_0$ 的 $n$ 次多项式（$n\geq1$），则 $\exists\alpha_1,\cdots,\alpha_n\in\mathbb{C}$，使得
\[
f(x)=a_0(x-\alpha_1)\cdots(x-\alpha_n).
\]
\end{prop}

\begin{proof}
对次数 $n$ 归纳。

$n=1$：$f(x)=a_0x+a_1=a_0\left(x+\dfrac{a_1}{a_0}\right)$（因 $a_0\neq0$），令 $\alpha_1=-\dfrac{a_1}{a_0}$ 即得。

$n=k+1$：由代数基本定理，存在 $\alpha\in\mathbb{C}$ 使 $f(\alpha)=0$；由引理 2.1，存在首项系数为 $a_0$、次数为 $k$ 的多项式 $g(x)$ 使
\[
f(x)=(x-\alpha)g(x)+f(\alpha)=(x-\alpha)g(x).
\]
对 $g(x)$ 用归纳假设，存在 $\alpha_1,\cdots,\alpha_k$ 使 $g(x)=a_0(x-\alpha_1)\cdots(x-\alpha_k)$，于是
\[
f(x)=a_0(x-\alpha)(x-\alpha_1)\cdots(x-\alpha_k). \qedhere
\]
\end{proof}

\begin{cor}
$f(x)$ 有且仅有 $n$ 个零点：$\beta\notin\{\alpha_1,\cdots,\alpha_n\}\Longrightarrow f(\beta)\neq0$。

记 $\beta_1,\cdots,\beta_s$ 是 $\alpha_1,\cdots,\alpha_n$ 中两两不同的数，则
\[
f(x)=a_0(x-\beta_1)^{e_1}\cdots(x-\beta_s)^{e_s},
\]
称 $e_i$ 是 $\beta_i$ 的重数。
\end{cor}

\begin{cor}[零点唯一性]
设 $f,g$ 分别是 $m,n$ 次多项式，$l=\max\{m,n\}$。若存在 $l+1$ 个数 $\alpha_1,\cdots,\alpha_{l+1}$ 使
\[
f(\alpha_i)=g(\alpha_i),\qquad i=1,\cdots,l+1,
\]
则 $f\equiv g$。
\end{cor}

\begin{proof}
令 $h=f-g$，则 $\deg h\le l$，且 $h$ 至少有 $l+1$ 个互不相同的零点。若 $\deg h\ge1$，由推论 2.3，非零的 $\deg h$ 次多项式至多有 $\deg h\le l$ 个零点，矛盾；若 $\deg h=0$，$h$ 为非零常数、没有零点，也矛盾。故 $\deg h=-\infty$，即 $h\equiv0$。
\end{proof}

\section{根与系数的关系}

记 $\sigma_i(\alpha_1,\cdots,\alpha_n)$ 表示从 $\alpha_1,\cdots,\alpha_n$ 中取 $i$ 个数相乘、再把所得 $\binom{n}{i}$ 个数相加（并约定 $\sigma_0=1$）。

\begin{lem}
\[
(x-\alpha_1)\cdots(x-\alpha_n)=\sum_{i=0}^{n}(-1)^i\sigma_i(\alpha_1,\cdots,\alpha_n)x^{n-i}.
\]
\end{lem}

\begin{proof}
对 $n$ 归纳，约定 $i>k$ 时 $\sigma_i(\alpha_1,\cdots,\alpha_k)=0$。

$n=1$：右边 $=\sigma_0x-\sigma_1=x-\alpha_1$。

$n=k+1$：
\[
\begin{aligned}
&\prod_{j=1}^{k+1}(x-\alpha_j)
=(x-\alpha_{k+1})\sum_{i=0}^{k}(-1)^i\sigma_i(\alpha_1,\cdots,\alpha_k)x^{k-i}\\
={}&\sum_{i=0}^{k}(-1)^i\sigma_i(\alpha_1,\cdots,\alpha_k)x^{k+1-i}
+\sum_{i=0}^{k}(-1)^{i+1}\sigma_i(\alpha_1,\cdots,\alpha_k)\alpha_{k+1}x^{k-i}\\
={}&x^{k+1}+\sum_{i=1}^{k}(-1)^i\Bigl[\sigma_i(\alpha_1,\cdots,\alpha_k)
+\sigma_{i-1}(\alpha_1,\cdots,\alpha_k)\alpha_{k+1}\Bigr]x^{k+1-i}
+(-1)^{k+1}\sigma_{k+1}(\alpha_1,\cdots,\alpha_{k+1})\\
={}&\sum_{i=0}^{k+1}(-1)^i\sigma_i(\alpha_1,\cdots,\alpha_{k+1})x^{k+1-i},
\end{aligned}
\]
其中用了 $\sigma_i(\alpha_1,\cdots,\alpha_{k+1})=\sigma_i(\alpha_1,\cdots,\alpha_k)+\sigma_{i-1}(\alpha_1,\cdots,\alpha_k)\alpha_{k+1}$，以及 $\sigma_{k+1}(\alpha_1,\cdots,\alpha_{k+1})=\sigma_k(\alpha_1,\cdots,\alpha_k)\alpha_{k+1}$。
\end{proof}

\begin{thm}[韦达定理]
设 $f(x)=a_0x^n+a_1x^{n-1}+\cdots+a_n=a_0\prod_{i=1}^{n}(x-\alpha_i)$，则
\[
a_i=a_0(-1)^i\sigma_i(\alpha_1,\cdots,\alpha_n),\qquad i=0,1,\cdots,n.
\]
特别地
\[
\alpha_1+\cdots+\alpha_n=-\frac{a_1}{a_0},\qquad
\alpha_1\alpha_2\cdots\alpha_n=(-1)^n\frac{a_n}{a_0}.
\]
\end{thm}

\begin{proof}
由引理 3.1，$a_0(x-\alpha_1)\cdots(x-\alpha_n)=\displaystyle\sum_{i=0}^{n}a_0(-1)^i\sigma_i(\alpha_1,\cdots,\alpha_n)x^{n-i}$，与 $f(x)=a_0x^n+a_1x^{n-1}+\cdots+a_n$ 逐项比较 $n+1$ 个系数即得。
\end{proof}

\section{实系数与整系数方程}

\begin{prop}[虚根成对]
设 $a_0,\cdots,a_n\in\mathbb{R}$，$a_0\neq0$。若 $\alpha\in\mathbb{C}$ 满足 $f(\alpha)=0$，则 $f(\overline{\alpha})=0$；且 $\alpha$ 与 $\overline{\alpha}$ 的重数相同。
\end{prop}

\begin{proof}
由 $a_0\alpha^n+a_1\alpha^{n-1}+\cdots+a_n=0$ 两边取共轭，注意 $\overline{a_i}=a_i$ 与共轭保持加法、乘法，得
\[
a_0\overline{\alpha}^{\,n}+a_1\overline{\alpha}^{\,n-1}+\cdots+a_n=0,
\]
即 $f(\overline{\alpha})=0$。故虚根 $\alpha\notin\mathbb{R}$ 必与 $\overline{\alpha}\neq\alpha$ 成对出现。

重数：记 $\overline{p}(x)$ 为把多项式 $p$ 的系数逐个取共轭所得的多项式，因 $f$ 的系数都是实数，故 $\overline{f}=f$。若 $(x-\alpha)^{e}\mid f$，即 $f(x)=(x-\alpha)^e h(x)$，两边系数取共轭得 $f(x)=(x-\overline{\alpha})^{e}\,\overline{h}(x)$，故 $(x-\overline{\alpha})^{e}\mid f$；把 $\alpha$ 与 $\overline{\alpha}$ 互换即得反向。
\end{proof}

\begin{thm}[整系数方程的因式定理]
设 $a_0,\cdots,a_n\in\mathbb{Z}$，$a_0\neq0$。若 $\alpha=\dfrac{q}{p}$（$p,q\in\mathbb{Z}$，$p>0$，$\gcd(p,q)=1$）是 $f$ 的零点，则
\[
p\mid a_0,\qquad q\mid a_n.
\]
\end{thm}

\begin{proof}
先证一个引理：若 $\gcd(p,q)=1$ 且 $p\mid qm$，则 $p\mid m$。事实上由辗转相除法（裴蜀等式），存在 $u,v\in\mathbb{Z}$ 使 $up+vq=1$，两边乘 $m$ 得 $upm+vqm=m$；因 $p\mid upm$ 且 $p\mid qm$，故 $p\mid vqm$，从而 $p\mid m$。

由 $f\!\left(\dfrac{q}{p}\right)=0$，
\[
a_0\frac{q^n}{p^n}+a_1\frac{q^{n-1}}{p^{n-1}}+\cdots+a_n=0,
\]
两边乘 $p^n$ 得
\[
a_0q^n+a_1q^{n-1}p+\cdots+a_{n-1}qp^{n-1}+a_np^n=0.
\]
于是 $a_0q^n=-\bigl(a_1q^{n-1}p+\cdots+a_np^n\bigr)$，右边每一项都含因子 $p$，故 $p\mid a_0q^n$；又 $\gcd(p,q)=1$ 推出 $\gcd(p,q^n)=1$，用上面的引理得 $p\mid a_0$。

同理由 $a_np^n=-\bigl(a_0q^n+a_1q^{n-1}p+\cdots+a_{n-1}qp^{n-1}\bigr)$，右边每一项都含因子 $q$，故 $q\mid a_np^n$；而 $\gcd(q,p)=1$ 推出 $\gcd(q,p^n)=1$，故 $q\mid a_n$。
\end{proof}

\end{document}
